% =====================================================================
%  R. Nielsen, AFTER DEATH, A CONTRIBUTION TO THE ELUCIDATION OF THE LAW OF SELFHOOD
%  In: For Idea and Reality.  A Journal, edited by R. Nielsen, Bj. Bjørnson and Rud. Schmidt.  Year 1871. Second Volume (Bind 5).
%  Copenhagen: Otto B. Wroblewsky, Publisher; Printed by G. S. Wibe.  Bind 5, pp. 193-223.  (Efter Døden, Bidrag til en Belysning af Selvhedsloven, in For Ide og Virkelighed.)
%  ENGLISH TRANSLATION, from transcription.tex (Danish), the text of record.
%  GENERATED by ebook-method/for-idee-1870-73/tr_build.py --write from this template and the translators' parts in
%  ebook-method/for-idee-1870-73/tr/ (TRANSLATOR-BRIEF.md; terms: GLOSSARY.md, its key terms approved by Hans
%  2026-10-13/14).  Conventions: TRANSLATION-PLAYBOOK.md.  Page numbers of the print in the margin.
% =====================================================================
\documentclass[12pt,a4paper,openany]{book}
\usepackage[utf8]{inputenc}
\usepackage[T1]{fontenc}
\usepackage{libertinus}
\usepackage{libertinust1math}
\usepackage{textalpha}
\usepackage[english]{babel}
\usepackage[protrusion=true,expansion=true]{microtype}
\usepackage{setspace}
\usepackage[margin=1.3in]{geometry}
\usepackage{fancyhdr}
\usepackage{hyperref}
\hypersetup{pdftitle={After Death}, pdfauthor={Rasmus Nielsen}, hidelinks}
\setlength{\headheight}{15pt}
\pagestyle{fancy}
\fancyhf{}
\fancyhead[LE,RO]{\thepage}
\fancyhead[C]{\textit{R. Nielsen: After Death (1871)}}
\setstretch{1.3}
\usepackage{marginnote}
\renewcommand{\thefootnote}{\arabic{footnote})}   % the print numbers each page's notes 1), 2)
\newcommand{\opage}[1]{\marginnote{\footnotesize\textit{[#1]}}}
\newcommand{\essayhead}[2]{\begin{center}{\large #1}\ifx&#2&\else\\[0.3em]{\normalsize #2}\fi\\[0.5em]\rule{3em}{0.4pt}\end{center}\medskip\noindent}
\newcommand{\signature}[1]{\par\smallskip\hfill #1\hspace*{2em}\par}

\begin{document}

% the title page of Bind 5 (1871-2-for-idee-og-virkelighed.pdf), in English
\begin{titlepage}
\begin{center}
\vspace*{3em}
{\Huge For Idea and Reality.}\\[1em]
\rule{8em}{0.4pt}\\[1em]
{\large A Journal,}\\[0.8em]
{\small edited}\\[0.4em]
{\footnotesize by}\\[0.6em]
{\large R. Nielsen,}\\[0.4em]
{\large Bj. Bjørnson and Rud. Schmidt.}\\[1em]
\rule{8em}{0.4pt}\\[0.6em]
{\large Year 1871. Second Volume.}\\[0.6em]
\rule{8em}{0.4pt}\\[4em]
Copenhagen.\\[0.4em]
Otto B. Wroblewsky, Publisher.\\[1.5em]
{\small Printed by G. S. Wibe.}\\[3em]
{\small Translated from the Danish}
\end{center}
\end{titlepage}

\essayhead{After Death,}{A Contribution to the Elucidation of the Law of Selfhood.}
% --- p. 193 ---
\setcounter{footnote}{0}%
\opage{193}How strange a guest Death is, after all! He moves about in the throng, yet minds his own business; he goes in and out
of men's houses like an old servant, a kind of official
person, with whose office the race has long since
become familiar. Mute and still, the man with the
scythe performs his accustomed work, an ordinary everyday business,
a work that all know, and that yet, when
the question is put concerning it, no one knows, for the work itself
is a riddle, an impenetrable secret. He
who has tried what sorrow is can join in the talk of sorrows,
and the memory of sorrows will make him eloquent; he
whom misfortune has once struck can join in the talk of misfortune,
and the reflection of terror and anguish can still lend flashes to his
speech; but he whom Death has struck cannot join
in the talk, cannot betray with a single word what has happened:
the dead man is silent about what he knows; the Bony One has,
as though for prudence's sake, pressed a seal of silence upon
his lips.

Just as Death's work is a secret, so also
the hour, the moment, in which it will present itself,
% --- p. 194 ---
\setcounter{footnote}{0}%
\opage{194}is on the whole a secret as well. He who on
official business travels about the country and inspects the tills
likes best to come upon the party concerned unawares, in order
to find out whether the accounts are in order. One might
almost believe that Death went on a similar errand, for it
loves to take men by surprise. Sometimes, through
sickness and pain, notice is given that now the moment is
near, now the Silent One is coming with his scythe; and then
it was only his shadow that went past the door: there is
no hurry; but one day, when all is well, and no one
is thinking of anything, he steps suddenly in and does his
duty. From this sensible people have drawn the lesson
that, in order not to be taken by surprise, we must always be prepared
for Death. It is a good and sound and practical
lesson, if only we knew in what the preparation is really to
consist.

To prepare oneself for Death is, in bourgeois practice,
commonly understood of what is called ``setting
one's house in order'', whether by amassing a fortune and making a
will, or by paying into the widows' fund or the annuity,
by, in short, securing for the survivors
a livelihood. When this has been duly seen to, a
bourgeois consciousness says confidently to itself: ``now I am prepared,
now I can die on any day it may be, for the account
is in order''. That such a preparation may be
in the highest degree respectable and necessary, we must concede;
it betrays prudent care and a loving mind to provide
for those nearest to one; a man who has fulfilled these conditions
can leave the world honorably and receive an
honorable burial: the funeral oration will not omit his praise.

But if preparing oneself for Death consists exclusively
% --- p. 195 ---
\setcounter{footnote}{0}%
\opage{195}in having the survivors provided for, then
there must surely be many who are exempt from preparing themselves;
the solitary have no need of it, for they have no one in the world
to provide for; nor have the rich, for the heirs will surely
know how to provide for themselves; the poor, yes, the poor might
sorely need a preparation, but they have, alas,
nothing with which to prepare themselves. To prepare oneself for
Death must, if the demand is to hold for all, necessarily
be understood of preparing oneself for a state after
death. But --- after death? this state is, after all, entirely
unknown; how can one prepare for
the entirely unknown? One can prepare for something
real, something determined beforehand, something of which one knows
that it will necessarily occur; one can, for example, be composed
and prepared for an earthquake, an eclipse of the sun,
an unavoidable war, the loss of one's fortune, etc. The
greater the certainty about what is to come, be it sad
or joyful, the more straightforward is the task
of preparation; but conversely, the more indefinite and indefinable
what is to come, the more painful and precarious is the
preparation. To prepare for the possible is
already far more difficult than to prepare for the
necessary; but to prepare for all possible possibilities,
and the entirely unknown looks like all possible
possibilities, is impossible. It is presumably this last consideration
above all that is to serve as an excuse for the many
who, filled with the bustle of life, have neither time nor
desire to brood over that enigmatic unknown ---
after death.

And what does Death say then? Well, Death says naturally
nothing; and yet, when the Hollow-eyed One draws near
% --- p. 196 ---
\setcounter{footnote}{0}%
\opage{196}and stares at one, it is lifelike as though Silence
itself became eloquent, as though it really said something.
``After death! Man, you lie to yourself when you
imagine that you know nothing at all and can know nothing
about the state after death. You call my nature enigmatic,
and enigmatic the land to which the journey leads;
for that very reason I advise you, before you go aboard my
emigrant ship, to make my closer acquaintance. Give
me a name, call me what you will, a crude force of nature,
an executioner of souls, a spirit from hell, an emissary
from heaven; but one thing is demanded unconditionally: you shall say
what you take me to be; you shall have a judgment, a decisive judgment
about me, you shall; for your judgment about me is
your own judgment about yourself: dig into your innermost, and there you will find
the key to my riddle''.

That man's judgment about Death is a judgment about his
own self can, with a little reflection, easily be made plain.
He who says that Death extinguishes the light of consciousness and
annihilates life says, in the same breath, that his
own self, the being on which his consciousness, his
feelings, representations, memories, etc. rest, vanishes
in death. It is the self that sees in the
stream of corruptibility its own image and its own law; it is
the self that plunges into the stream; it is the self that condemns
itself to death. ``But scientific physiology
has expressly proved that the law of selfhood is only
an organic natural law''. False, false! Physiology knows
nothing, nothing at all, about the law of selfhood. The laws to which
physiology refers all activities of life and soul
in the human body are mere physical laws,
laws for the selfless. The only thing physiology can
% --- p. 197 ---
\setcounter{footnote}{0}%
\opage{197}do for him who would ask its judgment on the question of immortality
is to confirm him in the delusion that the
soul, which cannot be thought without a body, can impossibly
form for itself a body of a higher order when the earthly one
is dissolved. But in order to be confirmed in this delusion,
one need not study physiology; one arrives more easily at the result
by quite simply surrendering to the testimony of the senses.
Homer is no scientific physiologist, but he has, as everyone with an open natural sense
must have, a strong, overwhelming impression that
the body is the soul's reality. The soul vanishes like
``a shadow, a dream'', when it leaves ``the whitish
bones''. Significant is the exclamation of surprise
with which the shade of Achilles at the entrance to the underworld
addresses Odysseus: ``Scion of Zeus, son of Laertes, Odysseus,
cunning one! Why have you ventured down to Hades, where
the dead have their home, beings without consciousness, phantoms of
men who have lived out their lives?'' That myth gives the shades
an existence in Hades is an obscure intimation of the
validity of the law of selfhood. At the very moment when the understanding, dazzled
by sense impressions, would make the soul into a nothing,
the immediate feeling of self reacts; imagination makes
the understanding's nothing into something again and holds the fleeting
something fast as a shadow.

The Greek conception of the realm of shadows is far from
having vanished with faith in Greek mythology;
the old myth lives still in thousands and thousands of souls,
frightens them in secret and fills them with hidden dread
of death; for the myth is in the obscurity of the mind, in the unclear
self, which is a stranger to its own being. What wonder
that a man in the midst of the splendor of the world of the senses has a
stifling foreboding of the realm of shadows, as long as his
% --- p. 198 ---
\setcounter{footnote}{0}%
\opage{198}real self is not yet much more than a shadow.
To live in such an intermediate state is intolerable; it is
almost like being buried alive.

In most people, accordingly, there is also an endeavor
to free themselves from the intermediate state; but the means
employed will not always answer to the end. There are
those who regard diversions and many kinds of outward-going
interests as a good narcotic, and there are those
for whom daily work, in combination with cares of livelihood
and all manner of temporal worries, gives quite enough to think
about; but all this is uncertain. An unexpected accident, a
death that really goes to one's heart, can suddenly darken
the vision and conjure the shades from Hades forth in the startled
imagination. Just as in the hot countries a snake or crocodile buried in the mire
can surface and frighten
the wayfarer, so the question, the uncanny
question of the state after death, can suddenly rise up
and throw mind and thoughts into turmoil. If the fear is to be overcome,
the matter must be taken up from the ground; this thought
seems indeed to be gaining more and more currency with advancing
enlightenment.

The modern knowing self-consciousness is already
tolerably clear with itself that this business of
the soul's immortality or, as we in this connection
prefer to express it, of the eternity of the self, is an old
fable, a nursery tale, which belongs in the nursery.
The self, an organic product of nature, arises in time and
perishes in time, this is one of the great discoveries
of metaphysical self-knowledge, the yield of centuries of
brooding, a fruitful result for clear, indeed
water-clear heads. If there is any deception in this, then it comes
% --- p. 199 ---
\setcounter{footnote}{0}%
\opage{199}not from science --- science is quite
innocent --- but from the self, which has misread and misused
science. It is not a physical natural law, but a
law of selfhood, that lies at the ground of the judgment. It is
the self that in this matter has to do exclusively with itself;
it is the self that, in order to get out of the stifling
uncertainty, with full composure prepares itself for death,
makes its will, and bequeaths to itself annihilation.

There are presumably those who regard the choice of such a way out
as a sign of despair, since the thought of an annihilation
is so dreadful. This is a rather naive view;
for he who is truly assured of his annihilation
has nothing to dread. People speak of the night of
death, of the darkness of the grave, of the loneliness in ``the cold,
narrow chamber'', etc., mere childish imaginings.
There is no darkness in the grave, no loneliness in the coffin;
for him who is annihilated in death there is, literally speaking,
neither light nor darkness, neither company nor loneliness:
the dead man himself is not even in the grave, for he
is not at all.

There are those who hold that everyone who has condemned
his soul to annihilation has committed a spiritual
suicide, a crime that in the end must bring with it
such weariness of life that the criminal ends by taking his own
life. In this way of looking at things, too, no small
misunderstanding lies hidden. Suicide is indeed always a
possible way out for him who, dead to eternity, lives
in his innermost being on a nothing; but it heightens self-esteem again to
know that man has his own life unconditionally in his own
hand. It is so far from the case that the thought of annihilation
% --- p. 200 ---
\setcounter{footnote}{0}%
\opage{200}should without further ado give individuals a distaste for living, that
the opposite is far rather the case. He who knows,
or at any rate believes he knows, within himself that all is over with
death can quite acquire a desire, an immense desire to use
and enjoy this brief present life. Once the thought of annihilation
has from the ground up severed the self from all that is
eternal and made it a temporal being, then
this temporal being must sink deep into the temporal, acquire
a greedy appetite for the temporal, nourish itself on it and
live in it as in its only proper, original element.
Only he knows how to live who rightly knows how to satisfy
life's demands. The highest satisfaction is
temporal self-satisfaction; temporal self-satisfaction
is enjoyment. To win the conditions of enjoyment, one must
work. Work is the means, enjoyment the end; work
is the law, enjoyment the gospel; that is the religion of the age.
The present is all, after death is nothing;
that is the philosophy of the age. The religion of the present day is wholly
absorbed in its philosophy: its faith is in truth one
with its knowledge.

But the ideas, yes, them we had nearly forgotten; the ideas
must of course be included, if the whole is to be complete.
The great world-ideas in politics and history, the ideas of beauty
in poetry and art, the ideas of the true in the sciences,
all this belongs essentially to the present
real life. The ideas have the unmistakable advantage
that they are at once both earthly and supra-earthly;
they furnish a guarantee that the omnipresent
eternal is in every way prepared to serve the temporal.
The modern eternity is liberal; it is a great convenience
about the eternal ideas that they are at bottom so modest.
% --- p. 201 ---
\setcounter{footnote}{0}%
\opage{201}The creditors of temporality are not modest; on the contrary,
they are importunate, crude realists, of whom one cannot
be rid until one has paid to the last farthing.
But the eternal ideas do not present themselves as creditors; they
are humane, disinterested beings. If one understands how to use
them, good! then they let themselves be used, by many even with
inconceivable success and astonishing, brilliant effect. If
one has no use for them and does not wish to have anything further to do
with them, good! then they withdraw with a smile,
and one notices no reproach. It shall then not
be said that the self which will be annihilated in death
has shut itself out from eternity, far from it! While other
souls must make do with hope, the hope that after
death they shall at last enter into eternity, this self
has made hope superfluous, and having forestalled Death, has,
instead of itself entering into eternity, let eternity
come to meet it by entering into time.

We have, it is said, full compensation for immortality
when we are willing to live for ideas. But in order to complete the law of selfhood,
we must surely take note of what it really means
to live for ideas. Naive souls, who readily let enthusiasm
run away with them, like to conceive the matter as though
living for an idea were the same as making oneself
the instrument of the idea, sacrificing everything for the idea's sake, having one's
life in the service of the idea. It is a beautiful turn of phrase, but
a skewed conception of the relation. On our
chosen standpoint it cannot possibly be the idea that is unconditionally to
rule over the self; according to the law of selfhood it must necessarily
be the self that is to rule the idea. Let us begin
with love; the enthusiasts say that they love
the idea for its own sake; that may be granted, inasmuch as
% --- p. 202 ---
\setcounter{footnote}{0}%
\opage{202}they mean thereby that it is the idea that procures them the
greatest pleasures, gives their life the richest sunshine,
and through all the vicissitudes of struggle and rest, sorrow and
joy affords them the highest self-satisfaction. One
can then love the idea for its own sake, as one can
love a young girl for her own sake: love
lasts as long as self-satisfaction lasts; self-satisfaction
lasts as long as desire lasts, and desire ---
as long as enjoyment lasts. But not a moment
longer. Take enjoyment away, and desire goes away too,
and when desire is gone, what is then to give love
nourishment? Duty! That is a self-deception. If there is no
disposition for eternity, no germ of eternal life, in the self, how
should it then be able to feel obliged to sacrifice itself
for something that is contrary to its nature? It is a monstrous
contradiction to want to oblige a temporal self to work
for eternal ends. The idea of justice, for example, no
human being with a clear understanding can ascribe unconditional
validity to; justice stands in the service of prudence
and is to be applied according to circumstances. The duty of justice
is a duty of compulsion, and every duty of compulsion ceases
with the compulsion. The weak must be just
toward the strong, but not conversely. The smaller state
ought to fulfill its obligations to the letter and show
itself just toward the greater; but does a
great power therefore always need to be just toward a
small power? The same holds of the idea of goodness. Everyone
who does the good must, if he is honest, admit
that he does not really do it for the good's sake, but
for his own sake, because it affords him satisfaction.
If, on the other hand, anyone would unconditionally do the good for
% --- p. 203 ---
\setcounter{footnote}{0}%
\opage{203}the good's sake, then he would have to be willing to do it even in the case that it were for himself an
evil, were repugnant to him; but that is a glaring
contradiction.

There is a philosophical ethics which, although it
quite sincerely denies the eternity of the self, the immortality of the soul,
nonetheless in a metaphysical preaching tone speaks with
unction of the unconditional claim of the good and the true upon
conscience, of man's obligation to do
the good even to the extreme degree of self-denial,
etc. It is an ambiguous ethics, which speaks with
two tongues. If the germ of eternity has been torn out of the self
to the last fiber, then there can be no judgment of eternity
in conscience. A temporal being can no
more be obliged to stake its life on an eternal
good than a fish can be obliged to sing or an
ape to study thoroughbass. Starting from the temporal, the
movement of the idea passes unceasingly through the eternal back into the
temporal; this circular movement is the exact
formula of the law of selfhood: the temporal is the point of departure, and the temporal
is the goal; eternity, the fleeting intermediary,
acquires significance only by casting luster upon the temporal
and furnishing stimulants for the present.

If one imagines that one can scare the life out of this
view by howling that the sacred is trampled
underfoot, that the moral powers are mocked, and the foundation
of human society is undermined, then a
Corybantic din against an ethical eclipse of the moon will be
without any significance. Here there is no ground for din,
but a serious summons to think. It by no means follows from the position set up that lawfulness, order
% --- p. 204 ---
\setcounter{footnote}{0}%
\opage{204}and morality should go under or even merely suffer
any real loss. Although the human
self is not eternal, it cannot nevertheless, as long as it lives,
cease to be rational. The rational thing is that man
in the consciousness of himself has a consciousness both
of the temporal and of the eternal, of the temporal as
reality, of the eternal as semblance. The truth is that the
eternal, which, when it is taken by itself and ascribed independent validity,
is an empty thought, an imagination, is nonetheless
of extraordinary effect when it is bent under the
temporal and goes together with the finite. It goes
with the eternal in consciousness as with space in outer
nature. Space in and for itself, pure infinite
space, is nothing, is only pure, infinite emptiness;
and yet real bodies could not exist
without space. So with the eternal. Eternity in
itself is nothing, and yet rational human life
cannot possibly exist without the thought of eternity. The law of gravity
shows how distances, that is, spatial relations,
gain influence on the activity of forces; the law of selfhood
shows how ideas, that is, relations of eternity, gain
influence on the activity of reason. Truth and
goodness, duty and right, these indispensable conditions of social life,
will not vanish from the earth because the
empty representation of the soul's immortality at last vanishes
for the self.

The foundation of all morality is sincerity, honesty
and uprightness. The great mass of industrious, staid,
decent citizens, who in all states are
the best supports of the existing order, give eloquent testimony
that the concepts: right and duty, are borrowed from the
% --- p. 205 ---
\setcounter{footnote}{0}%
\opage{205}temporal and find application only to the temporal.
Whatever higher intimations the individual may have of a true
holiness, a perfect justice, an eternal retribution,
whether he has these intimations from pastor and
parish clerk or from his own inner life, he must yet, if he
wishes to be sincere, admit that such ideal representations
are excluded from the exchange, from the marketplace, from the great world market,
that they do not hold at all in trade and conduct, that
they find, to speak plainly, no real application to
real things. ``Everyone is a thief in his trade''. If
all could agree to make this admission,
a number of high-flown phrases would indeed vanish,
but in return there would come far more
agreement between word and deed, far more
reliability, plainness and naturalness into social
relations.

To be sure, there is something revolting to the immediate feeling
in the fact that Death, in the moral sense too, makes
all equal. The noble man dies, and the scoundrel dies, and Death
draws a black line through the account, and erases the difference
between virtue and vice. This cool impartiality
makes, however, when we reflect aright, a grand moral
impression; it subdues the selfish and cools passion.

It is, after all, an absurdity that a man should
count as disinterested who here in the world does the good
without reward, and why? Because the person will not let himself
be satisfied with the trifle of reward that the present can offer him,
but screws the price up to such a height that
a whole eternity is needed to pay him his reward. What
beautiful disinterestedness! A poor wretched laborer, who
% --- p. 206 ---
\setcounter{footnote}{0}%
\opage{206}spurns a wage of ten rix-dollars on Monday in the
hope that by a strike it will rise to double,
if only he can wait till next Saturday, is not thus
disinterested. The much-praised disinterestedness that
sends the screwed-up bills in to the bureau of eternity
is a supreme self-interest, a cunning strike of
insatiable selfishness. Against this strike Death, the
strong employer, must hold his ground: we are all, after all, Death's
day laborers --- from Death's office no payment is made;
after death there is no wage, for after death there is
nothing.

From the other side too the thought, rightly viewed, is
uplifting, great and moral. When the unjust man ---
the robber, the deceiver, the murderer --- dies, then all
the upright feel how beneficent it is that Death keeps justice
and sets a limit to crime. Then the
upright kiss Death's hand and praise his good deeds.
But the exceedingly virtuous, those disinterested souls who
wait for the great reward, are not content with such a
simple act of justice. As a supplement to merit,
an extra gratuity to the reward yonder, they demand
into the bargain the refined enjoyment that, while they themselves are
rewarded throughout a whole eternity, the ungodly man shall in
their sight be punished throughout a whole eternity, shall without
mercy or pardon suffer the eternal torments of hell,
what an insatiable feeling of vengeance! It is, after all, a beautiful and
exalted trait in Death's character that he scorns to
give such a passion nourishment. Cold Death extinguishes
the fire of vengeance: after death there is no punishment, for
after death --- there is nothing.

In sharp opposition to the doctrine of annihilation stands
% --- p. 207 ---
\setcounter{footnote}{0}%
\opage{207}now the Gospel of the resurrection of the dead, a word of revelation
that is inseparably one with the Gospel of
Christ's resurrection. Either --- Or; either believe in
Christ's resurrection and reject the doctrine of annihilation, or
reject the Gospel of the Risen One and believe in annihilation;
a third is not given. To a superficial
view one would indeed think that there was
not merely a third, but a fourth, a fifth, a sixth, so
many exertions are made to avoid ``the
great extremes'' and to fit out for oneself a doctrine of immortality
more tempered to the line of thought of the age and of an enlightened reason.
But that leads to nothing; it \emph{is impossible}. The many
half-scientific projects of immortality that rationalistic
enlightenment has brought to light are in their
innermost ground not a hair better than the many unscientific
projects of perpetual motion that the half-educated
rogues of natural philosophy bring to light. Both
kinds of projects have this in common, that they stand in conflict
with the laws, the former with the laws of spirit, the latter with the
laws of nature.

The law of selfhood is the fundamental law of spirit. It is this
law to which every question about the state after death
must be led back; it is this law that must guide
us in judging the many parables, images
and symbols under which the beyond is described to us
in the divine word.

In the doctrine of annihilation the law of selfhood is, as has been
shown above, determined as a law of temporality; according to the Gospel
it becomes, as we shall now see, decisive to determine
as a law of eternity; one of two things: \emph{either} the nature of the soul
must be conceived in such a way that the self, sprung from the
% --- p. 208 ---
\setcounter{footnote}{0}%
\opage{208}the selfless, the material, after a temporal consciousness of
itself and of what is universal in reason returns to its
origin in the material and vanishes into the selfless;
or the soul must be conceived in such a way that it, sprung
from the divine self, and thus eternal according to its principle,
enters unconsciously into temporality, that on temporal
conditions it can develop the consciousness of the eternity
of its being, and, enriched with the yield of development, return
to its source. It will now appear that all that
the Gospel communicates to us with divine authority about
the resurrection of the dead is neither more nor less than
what can serve precisely to awaken a living view
and personal assurance of the fulfilment of the law of selfhood.

To understand this, it is necessary to lay aside
old prejudices and to sweep the musty, matted web of
theological-dogmatic representations out of consciousness.

When someone is about to travel to America
or to New Holland, it is in order that he should in every
possible way seek to procure the most reliable
information about the continent unknown to him, its
climate, its rivers and mountains, animals and plants, about the inhabitants,
what they live on, how they dress, and so
forth.  But in this way one cannot procure
information about the world to come, about the state
after death.  But, it is said, Christ and the Apostles
surely must know the facts; let us hold to their words,
and so we shall receive reliable information.  In this way
the theologians have from the beginning done their part to
fortify the congregations in the fixed idea that the coming
world was a kind of supraterrestrial Australia or
% --- p. 209 ---
\setcounter{footnote}{0}%
\opage{209}New Holland, from which, by the aid of passages of Scripture and visionary
considerations, treated, so to speak, in
a dogmatic-scientific manner, one was in a position to procure
the most reliable reports.  What
the case is with these reports can be seen on
every page of the history of dogma.

It has cost exertions to find out what
the case is with the ``intermediate state'', but it has
cost still greater exertion to get it quite exactly
settled how the resurrection of the
dead is actually to proceed.

``According to the concordant doctrine of the orthodox teachers of the church, the Scholastics and
the Protestant theologians, in the resurrection the same body is restored
with the same constituents as it had in earthly life,
with skin and hair, entrails, bones and teeth --- for how
else could gnashing of teeth take place in hell?
asked Jerome.  But now one was at once in embarrassment about the hairs.  Christ had said that not one hair
of our head should be lost (Luc. 21:18).  But
how hideous the body would become, exclaimed Augustine,
who connected this passage with the resurrection,
if the whole length of the head-hair and beard gradually
cut off, and perhaps also of the nails, were again
to grow back to it in the resurrection.  It was a stroke of luck
that the Lord had said in another place that all our
hairs were numbered (Matth. 10:30); for in this one saw a
comforting hint that with the restitution only
the number, not the length of the hairs, was meant.  But
what if someone had had too few hairs in this life?
Bodily defects of every kind, Augustine holds, are
% --- p. 210 ---
\setcounter{footnote}{0}%
\opage{210}corrected in the pious at the resurrection, excess of
fatness and leanness evened out; yet in such a way that
the basic features of the former shape, which are necessary
to make the person recognizable, are not effaced....
The imagination had great difficulty in picturing
how someone who, to use Luther's words, has
left a leg in Hungary, another whom wild beasts
have torn or dragged to death, or whom the fishes in the sea
have eaten, should at the resurrection again come by
these limbs\dots{}.  A more conspicuous difficulty
lay in the one case, which the Church Fathers
had already set up, that a man in a famine had
eaten of another's corpse: whether, that is, the piece of flesh in question
should rise as a part of the man to whom
it originally belonged, or of the one who had afterwards assimilated
it into himself by digestion?  In the latter case
the first lacked a definite, bounded piece of
his body, in the former the second lost all the constituents
he had appropriated by eating the piece.  Augustine
awards the disputed piece of flesh to the original possessor
and lets God give back to the other the
evaporated parts of flesh, after having gathered them from
the air.''\footnote[1]{D. F. Strauss.  The Christian Doctrine of Faith (Den kristelige Troeslære).  Vol. 2; 3rd section, \S\ 103.}

This only as a little foretaste for those whose teeth
water easily whenever they catch a whiff
from the theological kitchen.  A warming-up of the old
dishes one still finds in every second ``Christian dogmatics'',
only that the strength is considerably thinned out.  In the
% --- p. 211 ---
\setcounter{footnote}{0}%
\opage{211}older writers the eschatological investigations are conducted with great vividness
and naive boldness, in the more recent with anxious
caution and infinite insipidity: that is the whole difference.

We hold to the law of selfhood.  The soul is immortal,
for the self is eternal; this is no object of knowledge,
but an unavoidable object of faith.  ``I am the resurrection
and the life'', says Christ (Joh. 11:25--26): ``he
that believeth in me, though he were dead, yet shall he live, and
whosoever liveth and believeth in me shall never
die.  Believest thou this?'' --- Yes, believest thou this?
On the answering of this question every single
soul must then keep its whole attention gathered.  If
the question is misunderstood, the answer must be accordingly.
And it is in truth misunderstood when one
takes Christ's words as if he wished to say: ``You shall
trust blindly in me, for I am the outward, historically
personified authority.  When I tell you that there is
an eternal life, then there is an eternal life, and that must be
enough for you''.  With this misunderstanding the representation
of Christ's resurrection as an outward, directly provable
fact stands in close connection.  ``If we did not'',
it is said, ``have it in the letter of Scripture in black and white
that Christ is truly risen from the dead, then
we should know nothing, we should not know at all what hope we might
make for ourselves about our state after death; but now, since it
is shown by Scripture that Christ is risen, just as surely,
for historical reasons just as provably, as that Julius
Caesar was murdered by Brutus, and Napoleon deported
to St. Helena, now we know, yes now we know, that the
dead shall rise to eternal life''.

% --- p. 212 ---
\setcounter{footnote}{0}%
\opage{212}How much such an explanation conflicts with the law of selfhood
is obvious.  The self will see with its own
eye, the knowing self with the eye of knowledge, the believing self with the
eye of faith.  The assurance of an eternal life is in the deepest
sense a self-assurance; but an original self-assurance
cannot possibly be attained by forgetting oneself in order to
listen to another, looking away from oneself in order to look
at another, giving up one's own self and throwing oneself in the dust
before another.  Is this the spirit in Christ's words?  It is
surely a spiritless distortion.  Christ does not speak of himself
as of an outward object; he speaks of his
being.  He does not say: I have from heaven information
that after death there follows a life and a resurrection;
but he aims at the self and says to man:
I am what you inwardly seek, I am the resurrection
and the life.

Believest thou this?  ``No!'' answers modern criticism.
``I do not believe it, for I do not \emph{believe} at all.  I am, as
science can bear witness, not a believer, but a knower.
Faith is for the common people, empty doubting for the half-educated;
but men of thought with the clear concept win
firm conviction in knowledge''.  One doubtless does such
an utterance too much honor by crossing oneself at
its boldness, as at a ghastly sign of something
demonic; one comes nearer to the truth by taking
the declaration straightforwardly for what it is, a sign of
immaturity in thought, a complete misconception of
the law of selfhood.  The eternity of the self cannot, as has been shown elsewhere\footnote[1]{Cf. the author's Philosophy of Religion, pp.\ 28--46; 88--97; 390--93.}
possibly be cognized in objective knowledge, it must
% --- p. 213 ---
\setcounter{footnote}{0}%
\opage{213}necessarily be cognized in faith.  The thought of faith is in its kind
just as strictly convincing as the thought of knowledge.

The thought of faith grounded in the law of selfhood is in all
respects true to itself, has everywhere the same distinguishing mark.  The
weakened self must come to strength, the dim, clouded
self must clear itself in the light of the true self.
The soul's insight into its own being is conditioned by its insight
into the being of Christ.  The power of inwardness with which
the believer holds fast the hope of the resurrection rests on that
strength of feeling, that earnestness of will, thus on that power of
faith, of inwardness, with which he is able to grasp the
eternal life in Christ's resurrection.  The faith in immortality can
then be distorted in two ways, in that man either one-sidedly
relies on his own self-assurance or just as one-sidedly
builds on the objective contemplation of an outward historical
phenomenon.

Superficially viewed, it would indeed seem as if man,
according to the law of selfhood, were referred exclusively, or at any rate chiefly,
to himself alone.  The matter is then
conceived thus.  The individual says to himself: ``As
surely as I am immortal, I must by my own reason be
able to discover the germ of immortality in the
dispositions of my being.  I look at Christ; he had unshakable certainty
of the eternal life, that is true.  He foretold
his resurrection and went willingly to his death, there is something
uplifting in that, something wonderfully great and divine.
But all of that cannot help me, if I cannot
in myself, by myself attain the same certainty.
If I cannot attain it, what does it help then that I
read and read four disputed and highly discordant
accounts of the miracle of the resurrection?  If I can attain
% --- p. 214 ---
\setcounter{footnote}{0}%
\opage{214}it, then I must say, I suppose, that Christ has been a helper to me only in so far as, by his word and his example, he
has been to me an occasion for deeper self-understanding;
but this my self-understanding stands wholly
independent of the historical element in his resurrection''.

Has any man on these terms ever
grasped the eternal life?  It is no effeminate mysticism of piety
that we would recommend in opposition to this.  Christ
is not a selfish, ambitious ruler who will not
grant men a little ``moral independence'' for
fear that someone might perhaps be able to help himself
without needing him.  He has indeed souls enough
who, merely to be rid of the moral law, load him
with submissive attendance, fawn upon him and pay
his grace compliments.  For it is, as touched on
above, precisely the second breach of the law of selfhood that the soul
will merely behave passively toward the Revealer, will contemplate
in pious idleness his suffering, death and resurrection as
a sacred spectacle, an official gala performance, that objectively
unveils the secrets of the Eternal.  That is the Christianity of the enjoyment of ideas,
of aesthetic dogmatism and of phrases,
a quite agreeable Christianity, which has only the one fault,
that it never brings life and incorruptibility to light.

Against both errors of the age, moralism and aestheticizing
dogmatism, the law of selfhood holds firmly back.
The latter error would soon be eradicated, if
it did not have something opiate-like and fragrant, something
soothing, pleasant to an age nerve-weak in everything spiritual,
about it.  For everyone knows that it is an error,
and no one wishes to be openly known as holding it.  If
the error went about naked with death under
% --- p. 215 ---
\setcounter{footnote}{0}%
\opage{215}its arm, then even the children in the street would point fingers
and say: fie, what an ugly error that Christianity of
ease is.  But the error is not so ill-bred
as to go naked; it knows how to dress, and it
has costumes enough, more than any empress.

The error of moralism enjoys greater respect; there is
as it were something grand about it.  Yet anyone
who will look a little more closely can easily see that it bears heavily
a self-contradiction.  Yes, if the moralist could make
himself a Christ, could by his own power fill his
soul with truth, righteousness and love of God
and neighbour in such a measure that he could with good conscience
say to himself: ``I too am, before Abraham
was, the resurrection and the life'', then he would certainly be
right.  But to take the matter from this side surely never occurs to him;
he is, after all, a reasonable man.  So with him we
come upon another road.  He concedes
that his being is not unconditionally one with the eternal
truth, the eternal righteousness, the eternal love,
but that in the spirit of truth and love he feels himself
obliged to strive, strive and strive after righteousness.
Here is the fundamental contradiction; for what I merely strive
after, that I do not have; and what I do not have, but
yet strive after in the infinite, is precisely that which
I infinitely need.

So the moralist, the reasonable man, then needs
infinitely the assurance of the eternity of life, because he,
as the restlessness of his ethical striving shows, infinitely needs
eternal life.  On what terms does he reach the satisfaction,
the high goal toward which he strives?  Not in
% --- p. 216 ---
\setcounter{footnote}{0}%
\opage{216}and through himself; for in and through himself he has only the need.  Nor through a stranger, one standing outside;
for that which is to satisfy him in spirit and truth
must, even if it comes to meet him from without, yet
announce itself in him from within as the holy ground of life,
the bread of eternity that satisfies the soul because it is one with
the soul's original being.

This is the condition of faith, the fulfilment of the law of selfhood,
that the human self in infinite self-activity shall
draw its nourishment from Christ's self, shall receive life from his
life, light from his light, i.e.\ live by his resurrection.
The assurance of Christ's resurrection as a really historical
fact thus goes most inwardly together
with the personal appropriation of Christ's being.  But
this assurance is sacramental.  He who has not
taken the sacrament upon it that Christ is really risen
from the dead, he may insist on the historical element as much as
he likes, and grub physiologically in brains and nerves as deep as
he likes: he never gets the better of the smouldering doubt,
never comes to full assurance of eternal life after
death.

By looking away from the kernel and keeping to the shells,
i.e.\ by conceiving Christ's discourses not as a
divine clarification of the being of the self, a light of life and
of resurrection, which in the night of souls dispels the darkness
of death, but as a kind of statistical description of the state
in the other world, topographical contributions to a
travel map, so that the dead, like undaunted tourists, can
find roads and paths exactly everywhere; yes, what is more,
that the still living can beforehand almost
get to see all the sights both in heaven and hell,
% --- p. 217 ---
\setcounter{footnote}{0}%
\opage{217}Christian dogmaticians have made it easy for dissolving
criticism to stir the whole sum of heterogeneous
representations together into a self-contradictory confusion and to lower
the doctrine of the resurrection in the eyes of the so-called cultivated.
If the word of life about the resurrection of the dead is only to be
regarded as a doctrine, then our Lord Christ does not
even have the honor of being the first inventor of this doctrine;
one may think of the mythologies and, as regards his nearest
surroundings, especially of the Pharisees.  Christ has
thus merely improved the doctrine, added a couple of paragraphs to
the doctrine, and if it should perhaps, according to biblical-theological
researches, when these are continued thoroughly for yet
a couple of thousand years, finally turn out that the accounts of the gospels
are tolerably credible --- confirmed
the doctrine by his own resurrection.

A conspicuous testimony to the dogmatic
spiritlessness lies directly in this, that Christ's various
statements, when they are taken literally as factual information
about the beyond, can easily be interpreted so that
they come into manifest conflict with one another.  In the parable
e.g.\ of the rich man (Luk. 16) the transition
from death to life occurs suddenly, both to blessed joy and
to eternal torment; and an equally sudden transition is indicated
in the promise to the robber on the cross: (Luk. 23:43).  But
according to Joh. 5:28 the final awakening to life is
designated in such a way that the dead are first raised up from the graves when
they ``hear the voice of the Son of God''.  The departed must thus
rest in the bosom of the earth until the last day; for which reason
an old Arabian sect of Christians held that the soul
died with the body and was first at the final resurrection
called to life with the body.  Just as evident is,
% --- p. 218 ---
\setcounter{footnote}{0}%
\opage{218}according to criticism, the unclear wavering between a Jewish
carnality and an Alexandrian spirituality in the whole
description.  The carnality is that the state of the blessed
is likened to a great banquet.  ``Many shall come
from the east and the west and sit down with Abraham,
Isaac and Jacob in the kingdom of heaven'' (Matth. 8:11); for which reason
the congregation too believes in the resurrection of the flesh.
But the spirituality is that the carnal difference between
man and woman is effaced to such a degree that they become
pure angels (Matth. 22, Luk. 20, Mark. 12); an empty,
tedious spiritualism, which corresponds but little to the resurrection of
the flesh.

From this one will gather what good of a scientific kind
comes out of the endless wrestling between dogmatics
and criticism.  Everywhere that Christ speaks of the life to come he speaks in parables; but these parables are
pictures of life, in which the eternal truth itself takes shape.
As uplifting and soul-shaking, consoling and terrible
as these pictures of truth are when, in the light
of the spirit, they dawn upon the inner sight, fill the self and set
the powers of the mind in motion, so great again is the confusion
when either a selfishly excited imagination
takes the word in vain, or a sober, spiritually poor
understanding takes it into its head to set them up as prospects
for outward, sensory orientation.  Each parable, each
picture makes the truth vivid with such great peculiarity
and in such striking immediacy that here
there can be no thought of bringing the individual pictures
into objective systematic connection.  The connection
is in the fundamental view, not in the pictorial designation.
They are turning-points, great decisive principal points that
% --- p. 219 ---
\setcounter{footnote}{0}%
\opage{219}each time on a particular occasion, in a particular connection
are illuminated.

A couple of examples.  Christ says to his disciples
(Matth. 19:28): ``Verily I say unto you, that ye which
have followed me, in the regeneration, when the Son
of man shall sit in the throne of his glory, ye also shall
sit upon twelve thrones, judging the twelve tribes of Israel''.
But to the scribes and Pharisees, who wished to
see a sign from him, he says in another place (Matth.
12:41--42): ``The men of Nineveh shall rise up in the
judgment with this generation, and shall condemn it; for they
repented at the preaching of Jonas; and, behold, a greater
than Jonas is here.  The queen of the south shall rise up in the
judgment with this generation, and shall condemn it; for she
came from the uttermost parts of the earth to hear the wisdom of Solomon; and,
behold, a greater than Solomon is here''.  Who now will be in a
position to bring these two statements into a vivid connection,
when, mark well, they are to be understood literally
as descriptions of outwardly determinate occurrences taking place in a determinate
order?  Who is to judge first,
the Apostles or the Ninevites or the Queen of Sheba?
Not to speak of the fact that Christ himself, the Judge of all the world,
sits silent on his throne and leaves the judgment to
others.

These statements take on an entirely different significance
when we note the particular principal point which
supplies the motive in each of the great pictures.  It is,
as regards the last picture, the unbelief of the Pharisees that
occasions Christ to a threatening warning against the hardened
generation, to remind it what a terrible
judgment it is going to meet.  It shall in eternity be
% --- p. 220 ---
\setcounter{footnote}{0}%
\opage{220}revealed how high the Ninevites, who after all heard only one
Jonas, and the Queen of the South, who after all heard only one
Solomon, in reverence for the Holy, in sense of the Exalted,
stand above a generation that openly defies the truth
revealed in flesh and blood.  The judgment is at
once inward and outward; inward, for it is the generation's
own hardening by which it is judged;
outward, for the example of the Ninevites and the deep longing
of the Queen of the South are phenomena that furnish a measure
(Cf. Matth. 11:21--24).  Thus the law of selfhood is
fulfilled.  The judgment is in the self, in the inner life of each individual
man and thereby of the generation; but this judgment
in the inner life becomes the more terrible in the same degree as
the death-bearing consequences of blinding and hardening unfold
in the clear day of eternity, and the offence is measured with
the measure of reality.

In the statement to the disciples the principal point is entirely
different.  It is the well-known scene with the rich young man
and Christ's severe words about the sickness of riches and the rich, which
occasioned Peter's question: ``behold, we have forsaken
all and followed thee; what shall we have therefore?''
Full compensation! is the answer.  But this answer points
expressly to the final catastrophe of the ``regeneration'', and
receives life and color in the picture.  Full compensation! For
life, which in death lays aside only the corruptible, becomes
in eternity richer, fuller, more real than any
mortal can imagine, in the outward as in the inward.
Full compensation for privations and sufferings, for sighs
and pain, for scorn and contumely!  For the pressure and
burdens of finitude are essential conditions of the inwardness of understanding,
of the strength of self-development, of the life of love in
% --- p. 221 ---
\setcounter{footnote}{0}%
\opage{221}faith and hope.  Full compensation in power and authority, in
splendor and honor!  It is not a goal of self-interest, an
outward reward for outward deeds, that is held out here
in prospect; it is life itself, the eternal life, a gift so
rich that only the true followers, the souls who on
the condition of self-denial feel the law of selfhood as the law of love,
have in love itself the condition for
receiving the gift.  The gift is inward, the gift of the spirit,
life from Christ's life; but the inward takes shape in
the outward.  When the Christ-life unfolds its power and
reveals its glory: then they shall, but also only
they to whom honor is due, receive honor.

The distorted opposition between ``Jewish carnality''
and ``Alexandrian spirituality'' falls away of itself from this
point of view.  The life after death is, when
we hold to the spirit in Christ's words, not a life of bodiless
souls or natureless spirits; decisive weight is laid here
on a bodiliness of a higher order, on what we
can only cognize by cognizing an absolute limit
for human knowledge, a supernatural nature.  It
is the transformation of the corruptible nature to which Christ
refers in his answer to the Sadducees.  It is not
either the bodily or the psychical opposition
of man and woman that he wishes to efface; it is the Sadducees'
objections, built upon carnal representations, that
he rejects.  The human self is as little a
soul without a body as a body without a soul: it is in
its development to fully real peculiarity a psychical-bodily
being.  If now the original individuality of the human being
is not to be effaced but transfigured in the resurrection,
then the likeness to the angels cannot possibly be
% --- p. 222 ---
\setcounter{footnote}{0}%
\opage{222}understood of a fading of human individuality,
an effacement of the peculiar shape of man and woman; the passage
is, as the context expressly shows, to be understood of the purity and incorruptibility
in the personal eternal relation between man
and woman.  The angels are neither man nor woman;
an angel has not flesh and blood.  But Christ is man,
not woman; he who was transfigured in the resurrection has flesh
and blood.  Flesh and blood or, as the Risen One himself
has designated it, ``flesh and bones'', signifies in this context
not sensory, corruptible elements, but
the fully real human being: the congregation does not believe
in a carnal resurrection, but it believes in the resurrection
of the flesh.

\emph{After death --- comes judgment}; this is the earnestness of
the whole consideration.  But what is lacking in the present age
is spiritual powers to awaken the earnestness of the thought of death
in such a way that it becomes a really true and deep
earnestness.  The thundering zeal by which religious men of earnestness,
penitential preachers, who on account of the tickling shudder
it awakens in the awakened to hear of the outermost darkness,
of the pool of hell and the lake of brimstone, when the steam stands over
the devotion and the congregation is properly sulphured, can even
become fashionable preachers, has not always the testimony of the Spirit.
Sensible people think as follows: ``if we are not
angels, then neither are we devils.  Time brings counsel;
there is an intermediate state''.  Let one picture
to oneself a man, one of those who, following the direction of the age
in the religiously humane, have gradually prepared
themselves for annihilation, waking up after death and,
if we may speak so freely, poking his head into the High Hall
% --- p. 223 ---
\setcounter{footnote}{0}%
\opage{223}of Eternity! --- yes, he will be astonished, but will he therefore
at once be crushed as a sinner?  By no means!  As the
humane man looks about him and catches sight of the examining judge,
he makes a well-considered apology and says
straight out: ``I see that I have made a mistake and beg most earnestly for
forgiveness''.  A just and humane judge cannot possibly
condemn a soul that so modestly begs for forgiveness.
The Catholic Church takes purgatory seriously;
but it works only on an immature common people, for the doctrine
of purgatory is chemically outward.  The Protestant church theologians
reason up and down about the intermediate state; there
is an eternal damnation, and yet again it is no
eternal damnation.  It is supposed to be ``a dogmatic antinomy''
--- we have thoroughly studied the antinomy in ``the
Christian dogmatics'' --- but the antinomy is a self-made
speculation; in the antinomy there is ethicizing ``wit'',
but not religious earnestness.  If it is really to become
earnestness, then we must go back to the religious principle,
to the principle of selfhood, and follow the law of selfhood strictly in one
and all.  If there is in souls a need, an essential need of an
eternal life, then there must be deep grounds in the self.  The possibilities
of blessedness must be in there, yes! but also
the possibilities of perdition; for the one conditions the other.
To descend into the abysses of the self and raise the dim possibilities
up to the light of day is earnestness, the earnestness of life, the earnestness
of death: after death comes judgment.
\signature{R. \emph{Nielsen}.}

\end{document}
